\section{Introduction}
\label{sec:intro}

\subsection{The diagnosis bottleneck in modern operations}

A production microservice system in 2026 emits four kinds of telemetry continuously: \emph{structured logs} (every service writes JSON-formatted lines for every request it processes), \emph{distributed traces} (each user-visible action is followed across the chain of services that handled it, recorded as a tree of timed spans), \emph{metrics} (the rate of requests, the rate of errors, latency percentiles, container CPU and memory, all sampled every few seconds), and \emph{cluster events} (Kubernetes records every pod restart, every scheduling decision, every health-check failure). Together this stream amounts to tens of thousands of records per minute for a moderately busy cluster.

Most of this stream is uneventful. A small fraction matters. An even smaller fraction is so worth investigating that an engineer records it permanently in Jira --- the shared database of past incidents that becomes the team's institutional memory.

The hard part of an on-call engineer's job is no longer detecting that something has gone wrong: modern anomaly detectors (histogram baselines, learned classifiers on numeric features, statistical alert rules) are essentially solved for that question. The hard parts are three:

\begin{enumerate}
\item \textbf{Is this real, or is it noise?} A page can fire because of a real customer-impacting incident, a brief blip that self-recovered before anyone noticed, a side effect of a recent deployment that will settle on its own, or test traffic from a load-generator. Distinguishing these four categories before waking an engineer is itself a hard call.
\item \textbf{Have we seen this before?} The institution's memory of past incidents lives in Jira. Each past ticket carries a summary, a description, a timeline of investigation comments, a resolution, and links back to the telemetry that triggered it. If today's incident looks like one of those past tickets, the on-call engineer can read its resolution and act in seconds. Today this lookup is done manually --- the engineer guesses keywords and searches Jira --- which typically takes tens of minutes and depends entirely on the engineer remembering the right vocabulary.
\item \textbf{Is this something new?} If no past ticket matches, the failure mode might be genuinely novel and worth careful investigation, or the search might simply have failed. Confusing the two wastes engineering effort.
\end{enumerate}

The system we describe in this report answers all three questions simultaneously for every 5-minute telemetry window: a calibrated probability that the window is ticket-worthy, a ranked top-5 list of past Jira tickets that look similar, and a Boolean flag that fires when no past ticket plausibly matches.

\subsection{Why we had to build the dataset from scratch}

We considered using public datasets and concluded that none gave us what we needed.

The largest public log datasets (HDFS, BGL, Thunderbird) carry rich log streams but no ticket-side ground truth --- they tell you which lines were anomalous but never tell you which engineer filed which ticket for which incident. The largest public Jira dataset, TAWOS, contains 458{,}232 real Jira issues across 39 open-source projects (Moodle, MongoDB, Mesos, Mule, Hyperledger Fabric, and others), with a sophisticated split of free-text description versus pasted code blocks --- but the tickets are not paired with the telemetry that produced them. Real production telemetry from cooperating organizations is rarely shareable for legal and competitive reasons and, even when shareable, almost never reproducible: an external reviewer cannot regenerate the data to check a claim.

We built our own paired dataset under a controlled lab. The price is that the data is synthetic; the payoff is that every test number we report can be traced back to a specific fault-injection script, a specific telemetry export, a specific feature pipeline, and a specific git commit. We used TAWOS heavily, but only as a \emph{reference} for what real Jira tickets look like: their length distribution, the frequency of pasted code blocks, the engineering vocabulary, the distribution of resolution outcomes. TAWOS shaped how we humanized our synthetic tickets; it was never used as labeled training data.

\subsection{Why a hybrid model became necessary}

Over the course of the project we trained seven independent diagnostic pipelines, each representing a different bet about which signal in the telemetry matters most. A gradient-boosted tree classifier on numeric features won the triage task. A fine-tuned dense sentence encoder won top-1 retrieval. A hybrid sparse-plus-dense fusion won top-5 retrieval coverage. A small log-sequence transformer won on specific scenario families. A Cypher-traversal retriever over an LLM-built knowledge graph won at structural similarity. A three-stage LLM agent was the only pipeline that emitted a novelty flag at all.

Every column had a different winner. No single pipeline dominated. The natural next move would have been to train one bigger model and hope it would learn the union of these signals, but the gradient-boosted classifier had already saturated the triage metric at near-perfect performance (PR-AUC $= 0.9998$); replacing it with a transformer cannot beat $0.9998$ by more than a rounding error. Instead, we built a four-layer cascade that uses each pipeline for the sub-task it is best at, with explicit composition rules for how the layers feed each other. The cascade is what made the seven Pareto-incomparable pipelines compose into a single system that ties or beats every individual baseline on every metric.

\subsection{Headline result}

On 1{,}008 held-out test windows from the in-distribution split, the final cascade achieves Hit@5 $= 0.912$ (the gold past ticket is in the top-5 suggestions 91.2\% of the time), Hit@1 $= 0.722$ (the gold ticket is the very first pick 72.2\% of the time), Mean Reciprocal Rank $= 0.794$, strict PR-AUC $= 0.9998$, novel-precision $= 0.940$, and novel-recall $= 0.793$. The last number is the headline win of the project: it represents a $+388\%$ relative improvement over the previously locked baseline of $0.163$ at unchanged precision. The cascade is robust to memory noise (Hit@5 drops only $2\%$ at $50\%$ distractor ratio) and to family-distribution shift (novelty F1 within $11\%$ of in-distribution under leave-one-family-out cross-validation). It runs at $16$ microseconds per window at inference after the per-pipeline predictions are cached.

\subsection{Roadmap of this report}

Section~\ref{sec:dataset} explains how the dataset was built: what workload we instrumented, how we injected each kind of fault, what labels each window carries and what those labels mean, how we used the TAWOS public Jira corpus to ground our humanizer in real engineer voice, and how the distractor tickets are designed to fool the model. Section~\ref{sec:eval-setup} defines the metrics in plain language and explains the split rules. Section~\ref{sec:pipelines} explains each of the seven models in technical depth --- what the model architecture actually is, what library it comes from, what loss it is trained with, what hyperparameters were used, and what each model contributes to the cascade. Section~\ref{sec:cascade} specifies the four-layer cascade: how the layers feed each other, what the math is, and why each design choice was made. Section~\ref{sec:experiments} reports the experiments that validated the design: 13 baseline ablations followed by 8 targeted refinement experiments. Section~\ref{sec:results} consolidates the headline numbers and the robustness analyses. Section~\ref{sec:industrial} discusses cost, training time, inference latency, and what it would take to deploy this system in production. Section~\ref{sec:cross-app} addresses the most important question for external validity: how does the cascade generalize to a second microservice application with a different architecture, and what is our plan to answer that question for ICSE submission. Section~\ref{sec:limitations} states limitations and threats to validity, and Section~\ref{sec:conclusion} concludes.
